
% =============================================================
%  Chapter 1 — Introduction
% =============================================================
\chapter{Introduction}
\label{ch:introduction}

% ----- §1.1  Broad context: LLMs generating code -----
\section{LLM-Based Code Generation}
\label{sec:intro-context}

Large language models have advanced rapidly from completing individual
functions~\cite{chen2021humaneval} to resolving real-world software engineering
issues across entire repositories~\cite{jimenez2024swebench}.  Today, models
routinely generate multi-file backend applications that span frameworks,
databases, and authentication layers.  Benchmarks such as
HumanEval~\cite{chen2021humaneval} and SWE-bench~\cite{jimenez2024swebench}
have been instrumental in tracking this progress, focusing on functional
correctness at the function or pull-request level.

More recently, \textsc{BaxBench}~\cite{vero2025baxbench} introduced a
benchmark that evaluates LLM-generated \emph{complete} backend applications
across 392~tasks, 28~scenarios, 14~frameworks, and 6~programming languages.
Its results reveal that while the best-performing models achieve roughly 62\%
functional correctness, approximately half of the correct implementations
contain exploitable security vulnerabilities---highlighting that code
generation is far from a solved problem.  Complementary work by
\textsc{AutoBaxBuilder}~\cite{vonarx2025autobaxbuilder} automates the creation
of new benchmark scenarios and security exploits, enabling the benchmark to
scale and evolve alongside model capabilities.

These results establish that LLMs can produce functionally correct backend
code, albeit with significant remaining gaps in security.  However, both
benchmarks share a fundamental limitation in their evaluation methodology: they
execute each generated application in an isolated, single-machine Docker
container and measure only whether the code passes functional and security
tests.  No assessment is made of how the application would \emph{perform}
under realistic deployment conditions.


% ----- §1.2  Narrow the gap: performance under deployment -----
\section{Beyond Correctness: Performance Under Real Deployment}
\label{sec:intro-gap}

In production, backend applications do not run in isolation.  They are deployed
on container orchestration platforms such as
Kubernetes~\cite{burns2016borg,kubernetes}, where multiple replicas share
finite CPU, memory, and network resources across a cluster of worker nodes.
Achieving high throughput under these conditions requires careful
\emph{deployment configuration}---decisions about replica counts, CPU and
memory requests and limits, pod placement and affinity, database topology, and
connection pool sizing.

A functionally correct application can still fail to serve production traffic
if these parameters are poorly chosen.  Under-provisioned replicas lead to
request timeouts; over-provisioned ones waste cluster capacity.  Suboptimal
placement creates network bottlenecks.  Connection pool exhaustion causes
cascading failures under load.  Today, these configurations are tuned manually
by operations engineers or handled by reactive autoscaling mechanisms such as
the Horizontal Pod Autoscaler~(HPA)~\cite{kubernetes-hpa} and Vertical Pod
Autoscaler~(VPA)---neither of which reasons proactively about the interaction
between application behaviour and infrastructure layout.

Generating correct code is therefore a necessary but insufficient condition for
production readiness.  \emph{Deployment configuration is a critical and largely
unexamined dimension of LLM-based software generation.}  No existing benchmark
evaluates whether LLMs can reason about---and iteratively improve---the
deployment of the applications they generate.


% ----- §1.3  The gap stated explicitly -----
\section{The Missing Dimension: LLM-Driven Deployment Optimization}
\label{sec:intro-missing}

Concurrent work has demonstrated that LLM agents can perform iterative
performance optimization at other layers of the software stack.
\textsc{VibeServe}~\cite{kamahori2025vibeserve} uses a multi-agent loop to
synthesise bespoke LLM serving runtimes, achieving significant speedups over
general-purpose systems such as vLLM by specialising GPU kernels, memory
management, and scheduling logic to a specific model--hardware--workload
combination.  Similarly, Liu~et~al.~\cite{liu2025jitsystems} propose
\emph{Just-in-Time Systems}, in which an agent iteratively synthesises entire
key-value stores from specification cards, outperforming hand-tuned baselines
such as FASTER and RocksDB on all configurations tested.  At the code level,
\textsc{SWE-Perf}~\cite{he2025sweperf} benchmarks LLM capabilities for
repository-scale performance optimisation and finds a substantial gap between
model and expert performance.

These systems share a common structure: an iterative loop in which an agent
proposes a candidate, validates correctness, measures performance, and uses
structured feedback to guide the next iteration.  However, they all operate at
the \emph{implementation} level---optimising source code, data structures, or
system internals.  The \emph{deployment} level---how a fixed application is
configured, replicated, and placed on a cluster---remains unexplored.

This gap leaves a fundamental question unanswered:
%
\begin{quote}
  \emph{Can an LLM agent iteratively optimise the deployment configuration of
  a backend application on a Kubernetes cluster, guided by runtime performance
  feedback, to maximise sustainable throughput under fixed infrastructure
  constraints?}
\end{quote}


% ----- §1.4  This work -----
\section{This Work}
\label{sec:intro-thiswork}

In this thesis, we extend \textsc{BaxBench} with a Kubernetes-based iterative
benchmarking framework that addresses this question.  Given a
\textsc{BaxBench} scenario and a target Kubernetes cluster, the framework
orchestrates a multi-phase refinement loop:
%
\begin{enumerate}
  \item A \textbf{baseline phase} in which the LLM agent generates both
        application code and an initial deployment specification, which are
        validated, deployed, and benchmarked under adaptive load.
  \item An \textbf{iterative refinement phase} in which the agent receives
        structured performance feedback---load-test statistics, per-endpoint
        latency distributions, resource utilisation diagnostics, and the
        previous deployment's goodput---and decides whether to improve the
        application code or the deployment configuration before repeating the
        deploy--benchmark cycle.
\end{enumerate}
%
The deployment specification is expressed as a constrained YAML schema
(\texttt{spec.yaml}) that gives the agent control over replica counts, CPU and
memory requests and limits, pod placement via node affinity, database topology
(e.g.\ read replicas, connection pooling), and service-level parameters---while
the framework handles manifest rendering, Kubernetes resource management, and
orchestration.

Performance is measured using an adaptive load profile
(\texttt{k8s-goodput-plateau}) that automatically discovers the sustainable
throughput of each deployment.  Rather than fixing the request rate a priori,
the load controller ramps concurrent virtual users in steps until the
\emph{goodput}---the sustained rate of successful HTTP responses---stops
growing meaningfully.  This makes iteration scores directly comparable across
applications, deployment sizes, and refinement rounds without per-scenario
tuning.


% ----- §1.5  Contributions -----
\section{Contributions}
\label{sec:intro-contributions}

This thesis makes the following contributions:
%
\begin{enumerate}
  \item \textbf{An iterative benchmark framework} that extends
        \textsc{BaxBench} with Kubernetes deployment, adaptive load testing,
        and a multi-phase refinement loop
        (decision~$\to$ code~$\to$ spec~$\to$ deploy~$\to$ bench~$\to$
        feedback), enabling the evaluation of LLM-driven deployment
        optimisation alongside code generation.

  \item \textbf{A constrained deployment specification} (\texttt{spec.yaml})
        that exposes a structured search space for the LLM agent---covering
        replica count, resource requests and limits, database topology, and pod
        placement---while the framework enforces cluster-capacity validation
        and deterministic manifest rendering.

  \item \textbf{An adaptive load profile} (\texttt{k8s-goodput-plateau}) that
        automatically finds the sustainable throughput of any deployment
        without per-application rate tuning, using overload detection based on
        failure rate, tail latency, and goodput collapse thresholds.

  \item \textbf{An experimental evaluation} across multiple \textsc{BaxBench}
        scenarios, frameworks, and LLMs, characterising how iterative
        refinement affects deployment performance and examining the relative
        contribution of code changes versus deployment tuning.
        % TODO: Replace with concrete numbers once experiments are complete.
\end{enumerate}


% ----- §1.6  Thesis outline -----
\section{Thesis Outline}
\label{sec:intro-outline}

The remainder of this thesis is organised as follows.
%
\begin{description}
  \item[Chapter~\ref{ch:background}] provides background on
    \textsc{BaxBench}, Kubernetes, distributed load testing with Locust, and
    LLM-based code generation.  It also surveys related work on LLM-driven
    system optimisation and deployment automation.

  \item[Chapter~\ref{ch:methods}] describes the system architecture, the
    iterative experiment loop and its individual phases, the adaptive load
    profile and goodput metric, failure handling, and the constrained
    deployment specification.

  \item[Chapter~\ref{ch:setup}] details the experimental setup, including
    cluster configuration, selected models, scenarios, and hyperparameters.

  \item[Chapter~\ref{ch:results}] presents the experimental results: goodput
    trajectories across iterations, analysis of code versus deployment
    refinement, and representative case studies.

  \item[Chapter~\ref{ch:discussion}] discusses limitations, threats to
    validity, and comparison with alternative approaches.

  \item[Chapter~\ref{ch:conclusion}] summarises the findings and outlines
    directions for future work.
\end{description}
